\documentclass[12pt,a4paper]{article}
\usepackage[utf8]{inputenc}
\usepackage[T1]{fontenc}
\usepackage{amsmath,amssymb,amsfonts}
\usepackage{geometry}
\usepackage{booktabs}
\usepackage{hyperref}
\geometry{margin=2.5cm}

\title{The Boundary--Bulk Relation $\delta\rho/\rho = -2\psi$ from the Gauss--Codazzi Equations}
\author{Kevin Mu\~noz}
\date{}

\begin{document}
\maketitle

\begin{abstract}
The formalization of horizon-centered cosmological dynamics derives the relation $\delta\rho/\rho = -2\psi$ by linearizing the background Friedmann equation. While consistent, this is a verification rather than a derivation: it does not establish the geometric origin of the relation, does not specify its domain of validity, and does not provide the sub-horizon extension. This document closes that gap. We show that the Friedmann equation is the Gauss constraint for the FRW Cauchy hypersurface, and that its perturbation --- the full Hamiltonian constraint --- gives
\[
\frac{\delta\rho}{\rho} = -2\psi\!\left(1 + \frac{1}{3}\left(\frac{k}{aH}\right)^2\right)
\]
where $k$ is the comoving wavenumber. On super-horizon scales ($k \ll aH$), this reduces exactly to the Document~5 result. The sub-horizon correction is the Poisson equation, derived here as the Codazzi (momentum) constraint in the same formalism. The boundary--bulk relation is thus a rigorous consequence of the Einstein equations, valid in the separate-universe limit of the Gauss embedding structure.
\end{abstract}

\section{Motivation}

In Document~5 of the Causal Horizon Framework, the key perturbation relation
\begin{equation}
\frac{\delta\rho}{\rho} = -2\psi \tag{1}
\end{equation}
is derived in three lines: $R_A = 1/H$ gives $\delta H/H = -\psi$; the Friedmann equation $H^2 \propto \rho$ gives $\delta\rho/\rho = 2\delta H/H = -2\psi$. While correct, this derivation:

\begin{itemize}
  \item Uses the background Friedmann equation applied locally, without geometric justification
  \item Does not specify a gauge or domain of validity
  \item Does not show sub-horizon corrections
  \item Does not demonstrate that (1) is a consequence of the embedding geometry of the horizon in spacetime
\end{itemize}

The present document provides this derivation from the Gauss--Codazzi equations for the FRW Cauchy hypersurface.

\section{The Gauss--Codazzi Equations in 3+1 Form}

In the 3+1 (ADM) decomposition of spacetime, a Cauchy hypersurface $\Sigma_t$ (a surface of constant time $t$) is embedded in the 4D spacetime $(\mathcal{M}, g_{\mu\nu})$ with:

\begin{itemize}
  \item \textbf{Intrinsic metric} $h_{ij}$: the induced 3-metric on $\Sigma_t$
  \item \textbf{Extrinsic curvature} $K_{ij} = -\frac{1}{2}\mathcal{L}_n h_{ij}$: how $\Sigma_t$ bends in $\mathcal{M}$, with $n^\mu$ the future-pointing unit normal
\end{itemize}

The \textbf{Gauss equation} (contracted twice) gives the Hamiltonian constraint:
\begin{equation}
{}^{(3)}R + K^2 - K_{ij}K^{ij} = 16\pi G\,\rho_\text{ADM} \tag{2}
\end{equation}

where ${}^{(3)}R$ is the intrinsic Ricci scalar of $\Sigma_t$, $K = h^{ij}K_{ij}$ is the trace, and $\rho_\text{ADM} = T_{\mu\nu}n^\mu n^\nu$ is the energy density measured by the normal observer.

The \textbf{Codazzi equation} (contracted once) gives the momentum constraint:
\begin{equation}
D_j\!\left(K^{ij} - h^{ij}K\right) = 8\pi G\,J^i \tag{3}
\end{equation}

where $D_j$ is the covariant derivative on $\Sigma_t$ and $J^i = -T_{\mu\nu}n^\mu e^{i\nu}$ is the momentum current.

Together, (2) and (3) encode all information about how the matter content constrains the geometry of $\Sigma_t$ and its rate of expansion.

\section{The Friedmann Equation as the Background Gauss Constraint}

For a spatially flat ($k=0$) FRW background with metric
\begin{equation}
ds^2 = -dt^2 + a^2(t)\,\delta_{ij}\,dx^i\,dx^j \tag{4}
\end{equation}

the Cauchy hypersurface $\Sigma_t$ has:
\begin{equation}
{}^{(3)}R = 0, \qquad K_{ij} = -H\,h_{ij}, \qquad K = -3H, \qquad K_{ij}K^{ij} = 3H^2 \tag{5}
\end{equation}

Substituting into the Gauss constraint (2):
\begin{equation}
0 + 9H^2 - 3H^2 = 16\pi G\rho \tag{6}
\end{equation}

\begin{equation}
\boxed{H^2 = \frac{8\pi G}{3}\rho} \tag{7}
\end{equation}

The Friedmann equation is the Gauss constraint for the FRW Cauchy hypersurface. The apparent horizon radius $R_A = 1/H$ is thus determined by the embedding geometry of $\Sigma_t$ in spacetime.

\section{Perturbation of the Gauss Constraint}

We perturb around the FRW background in Newtonian (longitudinal) gauge:
\begin{equation}
ds^2 = -(1+2\Phi)\,dt^2 + a^2(t)(1-2\Phi)\,\delta_{ij}\,dx^i\,dx^j \tag{8}
\end{equation}

where $\Phi = \Psi$ (no anisotropic stress). The perturbed geometric quantities on $\Sigma_t$ are:
\begin{equation}
\delta\!\left({}^{(3)}R\right) = -\frac{4k^2\Phi}{a^2}, \qquad \delta K = 3(\dot\Phi + H\Phi), \qquad \delta(K_{ij}K^{ij}) = -6H(\dot\Phi + H\Phi) \tag{9}
\end{equation}

where $k$ is the comoving wavenumber. Linearizing the Gauss constraint (2) using background values (5):
\begin{equation}
\delta\!\left({}^{(3)}R\right) + 2K\,\delta K - \delta(K_{ij}K^{ij}) = 16\pi G\,\delta\rho \tag{10}
\end{equation}

Substituting:
\[
-\frac{4k^2\Phi}{a^2} + 2(-3H)\cdot 3(\dot\Phi+H\Phi) - (-6H(\dot\Phi+H\Phi)) = 16\pi G\,\delta\rho
\]
\[
-\frac{4k^2\Phi}{a^2} - 12H(\dot\Phi+H\Phi) = 16\pi G\,\delta\rho
\]

Dividing by $-4$, the linearized Gauss constraint takes the standard form:
\begin{equation}
\frac{k^2\Phi}{a^2} + 3H(\dot\Phi + H\Phi) = -4\pi G\,\delta\rho \tag{11}
\end{equation}

This is the \textbf{perturbed Hamiltonian constraint}, expressing the bulk energy density perturbation in terms of the geometry of the perturbed Cauchy hypersurface.

\section{Derivation of $\delta\rho/\rho = -2\psi$ (Super-Horizon Limit)}

\subsection{The growing mode from the Codazzi constraint}

The momentum constraint (Codazzi equation) in Newtonian gauge reads:
\begin{equation}
\dot\Phi + H\Phi = -4\pi G(\rho + P)\,V \tag{12}
\end{equation}

where $V$ is the velocity potential. On \textbf{super-horizon scales} ($k \to 0$), $V \propto k/aH \to 0$. For the growing mode, $\dot\Phi \approx 0$, i.e., $\Phi \approx \text{const}$.

\subsection{The Hamiltonian constraint in the long-wavelength limit}

Substituting $\dot\Phi = 0$ into (11) and taking $k \to 0$:
\begin{equation}
3H^2\Phi = -4\pi G\,\delta\rho \tag{13}
\end{equation}

Using the background Friedmann equation $4\pi G\rho = 3H^2/2$:
\begin{equation}
3H^2\Phi = -\frac{3H^2}{2}\,\frac{\delta\rho}{\rho} \tag{14}
\end{equation}

\begin{equation}
\frac{\delta\rho}{\rho} = -2\Phi \tag{15}
\end{equation}

In the super-horizon growing-mode limit in Newtonian gauge:
\begin{equation}
\psi \equiv -\frac{\delta H}{H} \approx \Phi \tag{16}
\end{equation}

Substituting:
\begin{equation}
\boxed{\frac{\delta\rho}{\rho} = -2\psi} \tag{17}
\end{equation}

This is the boundary--bulk relation of Document~5, now derived as the \textbf{super-horizon limit of the Gauss constraint}. It is valid for $k \ll aH$ and the growing-mode solution $\dot\Phi \approx 0$.

\section{Full $k$-Dependent Extension}

Retaining the $k^2\Phi/a^2$ term in equation (11) with $\dot\Phi = 0$:
\begin{equation}
\frac{k^2\Phi}{a^2} + 3H^2\Phi = -4\pi G\,\delta\rho = -\frac{3H^2}{2}\,\frac{\delta\rho}{\rho} \tag{18}
\end{equation}

\begin{equation}
\frac{\delta\rho}{\rho} = -\frac{2\Phi}{3H^2}\left(\frac{k^2}{a^2} + 3H^2\right) = -2\Phi\!\left(1 + \frac{k^2}{3a^2H^2}\right) \tag{19}
\end{equation}

Using $\psi \approx \Phi$:
\begin{equation}
\boxed{\frac{\delta\rho}{\rho} = -2\psi\!\left(1 + \frac{1}{3}\left(\frac{k}{aH}\right)^2\right)} \tag{20}
\end{equation}

This is the \textbf{full Gauss-constraint boundary--bulk relation}, valid for all $k$ in the growing-mode approximation.

\paragraph{Limiting regimes.}

\textbf{Super-horizon} ($k \ll aH$):
\begin{equation}
\frac{\delta\rho}{\rho} \approx -2\psi \tag{21}
\end{equation}

\textbf{Horizon crossing} ($k = aH$):
\begin{equation}
\frac{\delta\rho}{\rho} = -\frac{8}{3}\,\psi \tag{22}
\end{equation}

\textbf{Sub-horizon} ($k \gg aH$):
\begin{equation}
\frac{\delta\rho}{\rho} \approx -\frac{2}{3}\!\left(\frac{k}{aH}\right)^2\psi \tag{23}
\end{equation}
The Newtonian Poisson equation, with $\Phi = \psi$ playing the role of the gravitational potential.

\section{The Codazzi Equation as the Bulk Velocity Relation}

The momentum constraint (Codazzi equation, eq.~12) provides the complementary boundary--bulk relation for the velocity field:
\begin{equation}
V = -\frac{\dot\Phi + H\Phi}{4\pi G(\rho+P)} \tag{24}
\end{equation}

Using $4\pi G(\rho+P) = \varepsilon H^2$ and for the growing mode at super-horizon scales ($\dot\Phi = 0$, $k \to 0$):
\begin{equation}
V_\text{super-horiz} = -\frac{H\Phi}{\varepsilon H^2} = -\frac{\psi}{\varepsilon H} \tag{25}
\end{equation}

\section{Separate-Universe Interpretation}

The super-horizon limit (eq.~21) admits a clear physical interpretation in the \textbf{separate-universe approximation}: for $k \ll aH$, each Hubble-sized patch evolves as an independent FRW universe with its own local parameters $H_\text{loc}$ and $\rho_\text{loc}$.

Each patch satisfies the local Gauss constraint:
\begin{equation}
H_\text{loc}^2 = \frac{8\pi G}{3}\rho_\text{loc} \tag{26}
\end{equation}

Writing $H_\text{loc} = H + \delta H$ and $\rho_\text{loc} = \rho + \delta\rho$:
\begin{equation}
2H\,\delta H = \frac{8\pi G}{3}\,\delta\rho \tag{27}
\end{equation}

Using $4\pi G\rho = 3H^2/2$ and $\psi = -\delta H/H$:
\begin{equation}
\frac{\delta\rho}{\rho} = \frac{2\,\delta H}{H} = -2\psi \tag{28}
\end{equation}

This is identical to (17). The Gauss constraint derivation is the geometric justification for the separate-universe approximation: it holds precisely when $k \ll aH$.

\section{Gauge Specification and Validity}

Equation (17) is derived in \textbf{Newtonian gauge} with the identification $\psi \approx \Phi$. In a general gauge, the appropriate statement is:
\begin{equation}
\frac{\delta\rho}{\rho}\bigg|_\text{comoving} = -2\psi_\text{comoving} + \mathcal{O}\!\left(\frac{k^2}{a^2H^2}\right) \tag{29}
\end{equation}

The relation $\delta\rho/\rho = -2\psi$ holds precisely when:

\begin{enumerate}
  \item The super-horizon condition $k \ll aH$ is satisfied
  \item The growing-mode condition $\dot\Phi \approx 0$ holds
  \item The single-fluid adiabatic approximation applies
\end{enumerate}

For precision calculations, the full equation (20) should be used.

\section{Corrections to Document 5}

\paragraph{Equation (16) --- Now a theorem.} The relation $\delta\rho/\rho = -2\psi$ is upgraded from an algebraic consistency check to a theorem: it is the super-horizon limit of the linearized Gauss constraint (eq.~11), valid for $k \ll aH$.

\paragraph{Sub-horizon extension.} Document~5 did not provide the extension to finite $k$. The full result is equation (20).

\paragraph{Sign correction.} The Hamiltonian constraint originally had a sign error; the correct form is:
\begin{equation}
3H(\dot\Phi + H\Phi) + \frac{k^2\Phi}{a^2} = -4\pi G\,\delta\rho \tag{30}
\end{equation}

\section{Remaining Open Problems}

\begin{itemize}
  \item \textbf{Non-adiabatic perturbations.} For entropy perturbations ($\delta S \neq 0$), additional source terms appear in the Gauss constraint.

  \item \textbf{Beyond the growing-mode approximation.} A complete transfer function from $\psi$ to $\delta\rho/\rho$ for all $k$ and all times requires solving the full coupled perturbation equations.

  \item \textbf{2-surface Gauss--Codazzi for the apparent horizon.} A complementary derivation using the Gauss--Codazzi equations for the 2-surface (the apparent horizon itself) via the Raychaudhuri equation for null congruences would provide a more direct boundary-based justification for (17). This remains open.
\end{itemize}

\section{Summary}

\begin{table}[h]
\centering
\begin{tabular}{lll}
\toprule
Result & Equation & Origin \\
\midrule
Friedmann as Gauss constraint & $H^2 = 8\pi G\rho/3$ & eq.~(2) with FRW geometry \\
Perturbed Gauss constraint & $k^2\Phi/a^2 + 3H(\dot\Phi+H\Phi) = -4\pi G\delta\rho$ & eq.~(11) \\
Boundary--bulk relation (exact) & $\delta\rho/\rho = -2\psi[1 + (k/aH)^2/3]$ & eq.~(20) \\
Super-horizon limit & $\delta\rho/\rho = -2\psi$ & eq.~(17) \\
Sub-horizon limit (Poisson) & $\delta\rho/\rho = -(2/3)(k/aH)^2\psi$ & eq.~(23) \\
Velocity field (Codazzi) & $V = -\psi/(\varepsilon H)$ & eq.~(25) \\
\bottomrule
\end{tabular}
\end{table}

\section{Conclusion}

The relation $\delta\rho/\rho = -2\psi$ is the super-horizon, growing-mode limit of the linearized Gauss constraint for the FRW Cauchy hypersurface. The sub-horizon extension, given by equation (20), transitions to the Newtonian Poisson equation for $k \gg aH$. The Codazzi constraint provides the complementary relation for the bulk velocity field. These results complete the geometric justification for the boundary--bulk correspondence in Document~5.

\begin{thebibliography}{9}
\bibitem{adm2008} R. Arnowitt, S. Deser, and C. W. Misner, ``Republication of: The dynamics of general relativity,'' \textit{Gen. Relat. Gravit.} \textbf{40}, 1997 (2008).
\bibitem{bardeen1980} J. M. Bardeen, ``Gauge invariant cosmological perturbations,'' \textit{Phys. Rev. D} \textbf{22}, 1882 (1980).
\bibitem{hayward1996} S. A. Hayward, ``Gravitational energy in spherical symmetry,'' \textit{Phys. Rev. D} \textbf{53}, 1938 (1996).
\bibitem{cai2005} R.-G. Cai and S. P. Kim, ``First law of thermodynamics and Friedmann equations of Friedmann-Robertson-Walker universe,'' \textit{JHEP} \textbf{0502}, 050 (2005).
\bibitem{dodelson2003} S. Dodelson, \textit{Modern Cosmology}, Academic Press (2003).
\end{thebibliography}

\end{document}
